
\subsubsection{Perturbation-Based Experimental Design}

Standard benchmark evaluation provides single point estimates per dimension, yielding no information about correlations. We generate multiple observations by systematically varying intervention parameters (random seeds, learning rates, training data subsets) while holding the intervention type constant. Each replicate produces a performance change vector (ΔTruthfulness, ΔFairness, ΔRobustness), and correlation structure across replicates reveals dimension relationships.

\subsubsection{Intervention Method}

We use LoRA (Low-Rank Adaptation) with rank r = 8, scaling factor α = 16, applied to attention layers (c_attn modules in GPT-2). LoRA updates a small parameter subset (0.24\% of GPT-2's weights) while preserving pretrained knowledge. This provides a controlled intervention targeting specific architectural components.

\subsubsection{Datasets and Metrics}

Three trustworthiness dimensions were evaluated:
\begin{itemize}
\item \textbf{Truthfulness}: TruthfulQA multiple-choice (MC2), 817 questions measuring factual accuracy
\item \textbf{Fairness}: Bias Benchmark for QA (BBQ), 1,000 samples measuring stereotype bias
\item \textbf{Robustness}: Adversarial NLI Round 3 (ANLI R3), 1,200 samples measuring adversarial reasoning
\end{itemize}

Performance is measured as accuracy (percentage correct) on each benchmark. We compute performance deltas (Δ = post-intervention - baseline) and analyze pairwise Pearson correlations across replicates.

\subsubsection{Five-Hypothesis Validation Chain}

The study tests five hypotheses forming a mechanistic chain:

\textbf{H-E1 (Existence)}: Interventions targeting one dimension produce statistically significant correlations with other dimensions. Threshold: >80\% of dimension pairs show |ρ| > 0 with p < 0.01.

\textbf{H-M1 (Target Improvement)}: Parameter updates reliably improve the target dimension. Threshold: mean Δ > 0 with p < 0.05 and ≥70\% directional consistency across replicates.

\textbf{H-M2 (Representation Propagation)}: Parameter updates reshape internal representations across network layers. Measured via Centered Kernel Alignment (CKA) comparing pre- and post-intervention activations. Threshold: ≥50\% layers show detectable change or significant correlation with performance.

\textbf{H-M3 (Selective Coupling)}: Representation changes affect dimension pairs selectively rather than uniformly. Tests whether correlation patterns differ across dimension pairs using permutation tests. Threshold: at least one pair shows non-random correlation structure (p < 0.05).

\textbf{H-M4 (Architectural Replication)}: Correlation patterns replicate across model families. Tests directional consistency (positive, negative, or neutral correlation) across GPT-2, OPT, and Pythia. Threshold: ≥60\% replication rate for at least one dimension pair.

Each hypothesis has a gate condition determining whether to proceed. H-E1 and H-M1 use MUST_WORK gates requiring strong validation. H-M2, H-M3, and H-M4 use SHOULD_WORK gates allowing continuation with documented limitations.

\subsubsection{Models}

Three transformer architectures were tested:
\begin{itemize}
\item GPT-2 (124M parameters): baseline transformer architecture
\item OPT-350M (350M parameters): Meta's transformer variant
\item Pythia-410M (410M parameters): EleutherAI's transformer trained on the Pile
\end{itemize}

Model selection prioritized architectural diversity within the transformer family while maintaining computational feasibility.
